\documentclass[10pt,journal,compsoc]{IEEEtran}
\usepackage{cite}
\usepackage{amsmath,amssymb,amsfonts}
\usepackage{algorithmic}
\usepackage{algorithm}
\usepackage{graphicx}
\usepackage{textcomp}
\usepackage{xcolor}
\usepackage{booktabs}
\usepackage{multirow}
\usepackage{array}
\usepackage{url}
\usepackage{hyperref}

\hypersetup{
    colorlinks=true,
    linkcolor=blue,
    citecolor=blue,
    urlcolor=blue
}

\def\BibTeX{{\rm B\kern-.05em{\sc i\kern-.025em b}\kern-.08em
    T\kern-.1667em\lower.7ex\hbox{E}\kern-.125emX}}

\begin{document}

\title{ScopeGuard: Identity-Scoped Execution Control and Dynamic Permission Authorization for AI Agents in Cloud Infrastructure}

\author{Binod~Prasad~Joshi,
        Cloud~Infrastructure~Security~Research~Group%
\IEEEcompsocitemizethanks{\IEEEcompsocthanksitem B. P. Joshi is with the Department of Computer Science and Cloud Infrastructure Security Research Group.\protect\\
E-mail: research@scopeguard.io}%
\thanks{Manuscript received August 4, 2026; revised August 4, 2026.}}

\markboth{IEEE Transactions on Dependable and Secure Computing,~Vol.~21, No.~4,~August~2026}%
{Joshi \MakeLowercase{\textit{et al.}}: ScopeGuard: Identity-Scoped Execution Control for AI Agents}

\IEEEtitleabstractindextext{%
\begin{abstract}
The rapid integration of Large Language Model (LLM)-powered autonomous agents into cloud management workflows promises to automate complex infrastructure-as-code (IaC) provisioning, resource scaling, and operational monitoring. However, deploying off-the-shelf AI agents with static cloud credentials introduces severe security vulnerabilities: agents granted broad Identity and Access Management (IAM) permissions can exhibit unintended scope creep, succumb to indirect prompt injection attacks, or execute destructive infrastructure actions without deterministic identity attribution. Traditional access control models—such as static Role-Based Access Control (RBAC) and declarative Policy-as-Code engines (e.g., OPA, AWS Cedar)—fail to dynamically infer permission boundaries from natural language task objectives prior to execution. We present \textbf{ScopeGuard}, an asynchronous, identity-scoped proxy framework designed for real-time permission authorization and execution control of AI agents in cloud infrastructure. ScopeGuard intercepts outbound agent tool calls, attaches cryptographically verifiable identity metadata (Agent ID, Session UUID, declared task, and reasoning chain), and dynamically compiles a task-aware permission scope using a structured cloud action taxonomy. Intercepted API calls are evaluated against this computed scope under a fail-closed policy engine supporting three ablation operational modes (\textit{passthrough}, \textit{tagging\_only}, and \textit{full}). We empirically evaluate ScopeGuard across a benchmark suite of 200 infrastructure management tasks (100 legitimate operations and 100 adversarial/off-scope attacks across 10 security threat categories). In \textit{full} enforcement mode, ScopeGuard achieves an overall accuracy of \textbf{89.0\%}, a scoping precision of \textbf{97.8\%}, an adversarial block rate of \textbf{88.0\%}, and a false positive rate of \textbf{2.0\%} on legitimate tasks, while introducing a minimal average latency overhead of \textbf{108.9 ms}—well within the 200 ms real-time inference budget. We validate live execution against real Amazon Web Services (AWS) infrastructure endpoints including Amazon RDS, EC2, S3, and CloudWatch.
\end{abstract}

\begin{IEEEkeywords}
AI Agent Security, Infrastructure Access Control, Cloud Governance, Prompt Injection Defense, Principle of Least Privilege, Real-Time Authorization Proxy, AWS Identity Scoping, Empirical Security Evaluation.
\end{IEEEkeywords}}

\maketitle
\IEEEdisplaynontitleabstractindextext
\IEEEpeerreviewmaketitle

\section{Introduction}\label{sec:intro}
\IEEEPARstart{C}{loud} computing infrastructure forms the backbone of modern enterprise computing, underpinning microservice architectures, web search, database storage, and AI inference workloads \cite{ayers2018}, \cite{jiang2026}. Operating these datacenters and public cloud environments requires executing multi-step infrastructure management tasks via Cloud API endpoints (e.g., AWS SDK \texttt{boto3}), Infrastructure-as-Code (IaC) CLI tools (e.g., Terraform), and deployment pipelines (e.g., GitHub Actions). As infrastructure grows in scale, heterogeneity, and operational complexity, human operations teams face severe bottlenecks in manual resource provisioning, security auditing, and incident remediation \cite{hadary2020}.

Recent breakthroughs in agentic AI offer a promising path forward. Autonomous AI agents—LLM-powered software systems that reason about goals, generate execution plans, and invoke external tools \cite{schick2023}, \cite{yao2023}—are rapidly being adopted to automate database scaling, cluster management, and CI/CD workflow execution \cite{hong2024}, \cite{yang2024}. Rather than replacing human operators, agentic frameworks amplify engineering productivity by executing complex operational sequences across cloud control planes.

However, simply granting autonomous agents direct access to cloud control planes creates critical security vulnerabilities:
\begin{enumerate}
    \item \textbf{Lack of Dynamic Task Scoping (\textit{Not Least-Privileged})}: Standard cloud IAM assigns static roles to service principals or container environments. To enable an agent to perform diverse operational tasks, administrators frequently assign broad permissions (e.g., \texttt{AdministratorAccess} or wildcard \texttt{s3:*}). Consequently, an agent instructed to "check database status" retains permission to terminate production EC2 instances or modify IAM security groups.
    \item \textbf{Vulnerability to Indirect Prompt Injection (\textit{Not Protected})}: Autonomous agents parse untrusted external content (e.g., application log traces, git commits, or external API responses). Adversaries can embed malicious instructions within these data sources, hijacking the agent's plan and forcing it to execute unauthorized, off-scope operations \cite{shavit2023}.
    \item \textbf{Absence of Session-Level Identity Attribution (\textit{Not Attributable})}: Cloud API audit logs (e.g., AWS CloudTrail) record only the static IAM role assumed by the agent container. They fail to bind individual API requests to specific agent instances, session UUIDs, declared user objectives, or underlying LLM reasoning steps, creating severe forensic gaps during security incident response.
\end{enumerate}

\begin{table*}[t]
\caption{Systemic Comparison of Access Control Paradigms for Autonomous AI Agents in Cloud Infrastructure}
\label{tab:systemic_comparison}
\centering
\begin{tabular}{lccccc}
\toprule
\textbf{Control Paradigm} & \textbf{Dynamic Task Scoping} & \textbf{Runtime Identity Attribution} & \textbf{Sub-200ms Overhead} & \textbf{Cloud API Awareness} & \textbf{Active Pre-Execution Block} \\
\midrule
AWS IAM / Static RBAC & \textbf{No} (Static Roles) & \textbf{No} (Role-level only) & \textbf{Yes} (Native) & \textbf{Yes} (Native AWS) & \textbf{No} \\
Open Policy Agent (OPA) & \textbf{No} (Hardcoded Rego) & \textbf{Partial} (Token-based) & \textbf{Yes} ($<$10ms) & \textbf{Partial} (Generic JSON) & \textbf{Yes} \\
AWS Cedar & \textbf{No} (Hardcoded Cedar) & \textbf{Partial} (Principal-based) & \textbf{Yes} ($<$10ms) & \textbf{Yes} (Native AWS) & \textbf{Yes} \\
NVIDIA NeMo Guardrails & \textbf{No} (Text Content Only) & \textbf{Partial} (Session-level) & \textbf{No} ($>$300ms) & \textbf{No} (Generic Text) & \textbf{Partial} (Text Layer) \\
RIVA (2026) & \textbf{No} (Post-hoc Drift) & \textbf{No} & \textbf{N/A} & \textbf{Yes} & \textbf{No} (Post-execution) \\
\textbf{ScopeGuard (Ours)} & \textbf{YES (Task-Aware)} & \textbf{YES (UUID4 + Task + Chain)} & \textbf{YES (108.9ms Avg)} & \textbf{YES (AWS/Terraform)} & \textbf{YES (Fail-Closed Proxy)} \\
\bottomrule
\end{tabular}
\end{table*}

\subsection{Our Work: ScopeGuard}
We present \textbf{ScopeGuard}, an asynchronous, identity-scoped policy proxy and authorization engine positioned between AI agent tool-use interfaces and cloud infrastructure backends. ScopeGuard enforces the Principle of Least Privilege dynamically at runtime: when an agent submits a tool call, it must supply its declared task objective. ScopeGuard intercepts the call, injects session identity metadata, compiles the allowable permission set for the declared task, and evaluates the call against a configurable policy engine before forwarding it to cloud endpoints.

This paper makes the following primary contributions:
\begin{itemize}
    \item \textbf{Formal Architectural Framework}: We design a lightweight, asynchronous proxy architecture combining identity tagging, secret redaction, rule-based permission compilation, fail-closed enforcement, and persistent SQLite audit logging.
    \item \textbf{Dynamic Task-to-Scope Compiler}: We construct a deterministic classifier operating over a structured cloud action taxonomy (\texttt{proxy/schemas/permissions.json}) that maps natural language task prompts to explicit API permission boundaries with implicit read-grant inheritance.
    \item \textbf{Empirical 600-Call Ablation Study}: We design a 200-task evaluation benchmark (100 legitimate operations and 100 adversarial attacks across 10 security threat categories) and evaluate performance across three operational modes (\textit{passthrough}, \textit{tagging\_only}, and \textit{full}).
    \item \textbf{Real AWS Validation}: We demonstrate live execution against real Amazon Web Services (AWS) endpoints (Amazon RDS, EC2, S3, CloudWatch) in \texttt{us-east-1}, showing that ScopeGuard achieves an \textbf{88.0\% adversarial block rate} and \textbf{97.8\% scoping precision} with only \textbf{108.9 ms latency overhead}.
\end{itemize}

\section{Background and Threat Model}\label{sec:background}

\subsection{Cloud Control Planes and IAM Limits}
Cloud control planes manage compute (EC2), storage (S3), relational databases (RDS), IAM security policies, monitoring (CloudWatch), and serverless execution (Lambda). Authorization is governed by JSON policy documents defining \texttt{Effect}, \texttt{Action}, \texttt{Resource}, and \texttt{Condition}.

Static IAM policies assume predictable microservice workloads. When assigned to autonomous AI agents, static policies present an insurmountable dilemma: over-provisioning grants broad administrative rights exposed to prompt injection, whereas under-provisioning causes agent execution failures whenever an agent chooses an alternative valid tool path.

\subsection{Policy-as-Code & Content Guardrails Limitations}
Policy-as-Code engines (OPA, AWS Cedar) evaluate declarative rules over JSON request structures. However, they require human security teams to hand-craft static rules for every resource ARN and API action. They lack the semantic capacity to parse an agent's natural-language task prompt (e.g., inferring that "Check DB health" allows \texttt{rds:DescribeDBInstances} but prohibits \texttt{iam:DeleteRole}).

Content guardrails (NeMo Guardrails, Lakera Guard) filter text generations at the prompt level. However, they operate outside the structured API call execution path and cannot inspect JSON parameter payloads sent to \texttt{boto3} or Terraform CLI.

\subsection{Threat Model and Security Scope}
We define a threat model comprising six primary attack vectors against cloud-connected AI agents:
\begin{enumerate}
    \item \textbf{Scope Creep (Off-Scope Action)}: Agent attempts API operations unrelated to its stated objective (e.g., deleting buckets during a log-viewing task).
    \item \textbf{Direct Prompt Injection}: Adversary manipulates the user prompt to force unauthorized execution.
    \item \textbf{Indirect Log/Input Injection}: Adversary embeds malicious payload in application logs parsed by the agent.
    \item \textbf{Privilege Escalation}: Agent attempts to attach administrative IAM policies (\texttt{iam:PutRolePolicy}).
    \item \textbf{Resource Destruction}: Agent attempts irreversible infrastructure deletion (\texttt{ec2:TerminateInstances}, \texttt{rds:DeleteDBInstance}).
    \item \textbf{Unauthorized Ingress Expansion}: Agent opens public security group ports (\texttt{ec2:AuthorizeSecurityGroupIngress} to \texttt{0.0.0.0/0}).
\end{enumerate}

\section{ScopeGuard System Architecture}\label{sec:architecture}

ScopeGuard functions as an asynchronous reverse proxy server implemented in FastAPI, intercepting requests submitted to the \texttt{POST /execute-tool} endpoint.

\subsection{Component Breakdown}

\subsubsection{Identity Tagger (\texttt{proxy/identity.py})}
Every incoming tool call request is assigned a unique UUID4 tracking identifier. The Identity Tagger binds:
$$\text{Metadata} = \langle \text{AgentID}, \text{SessionID}, \text{UUID4}, T_{\text{decl}}, \text{Timestamp}, \text{ReasoningExcerpt} \rangle$$
Before logging parameter payloads, recursive regular expression matching (\texttt{\_SECRET\_PATTERNS}) redacts sensitive credentials (passwords, AWS secret access keys, JWT tokens) with \texttt{[REDACTED]} placeholders (NFR5).

\subsubsection{Task-to-Scope Compiler (\texttt{proxy/classifier.py})}
The compiler parses the declared task string $T_{\text{decl}}$ against an extensible JSON taxonomy schema (\texttt{proxy/schemas/permissions.json}). It tokenizes $T_{\text{decl}}$, matches action keywords across resource modules (S3, EC2, IAM, RDS, CloudWatch, Lambda), and extracts allowed actions $S_{\text{allowed}}$ and risk levels $R \in \{\text{LOW}, \text{MEDIUM}, \text{HIGH}, \text{CRITICAL}\}$.
If implicit read grants are enabled, permitted write actions automatically inherit corresponding read actions (e.g., \texttt{s3:CreateBucket} grants \texttt{s3:ListBuckets}).

\subsubsection{Policy Engine (\texttt{proxy/policy.py})}
Evaluates requested tool action $a_{\text{req}}$ against $S_{\text{allowed}}$ under three ablation operational modes:
\begin{itemize}
    \item \texttt{passthrough}: Unfiltered baseline; logs request and returns \texttt{ALLOWED}.
    \item \texttt{tagging\_only}: Attaches identity metadata and audit logs without blocking.
    \item \texttt{full}: Active enforcement. If $a_{\text{req}} \in S_{\text{allowed}}$, decision is \texttt{ALLOWED}. If $a_{\text{req}} \notin S_{\text{allowed}}$, decision is \texttt{BLOCKED}.
\end{itemize}
Operates under a fail-closed paradigm (\texttt{SCOPEGUARD\_FAIL\_CLOSED=True}): internal execution errors default to \texttt{BLOCKED} (NFR2).

\subsubsection{SQLite Audit Logger (\texttt{proxy/audit.py})}
Asynchronously persists structured \texttt{AuditEntry} records in SQLite (\texttt{data/scopeguard\_audit.db}) with indexed querying over Agent ID, Session ID, decision, and risk level.

\subsubsection{Backend Executor (\texttt{proxy/executor.py})}
Dispatches authorized calls to LocalStack Docker emulation (port 4566) or live AWS Cloud APIs via \texttt{boto3}.

\section{Formalization and Algorithmic Design}\label{sec:formalization}

Let $\mathcal{A}$ denote the set of cloud API actions in the taxonomy. An incoming tool request $Q$ is represented as:
\begin{equation}
Q = \langle \text{agent\_id}, \text{session\_id}, T_{\text{decl}}, a_{\text{req}}, P \rangle
\end{equation}
where $T_{\text{decl}}$ is the declared task string, $a_{\text{req}} \in \mathcal{A}$ is the requested action, and $P$ is the parameter payload.

\begin{algorithm}[t]
\caption{Task-to-Scope Permission Compilation}
\label{alg:compiler}
\begin{algorithmic}[1]
\REQUIRE Declared task $T_{\text{decl}}$, Action Taxonomy Schema $S$
\ENSURE Allowed Action Set $S_{\text{allowed}}$, Max Risk $R$
\STATE $Tokens \leftarrow \text{TokenizeAndStem}(T_{\text{decl}})$
\STATE $S_{\text{allowed}} \leftarrow \emptyset$, $R \leftarrow \text{LOW}$
\FOR{each $res\_type \in S.\text{resource\_types}$}
    \FOR{each $action, meta \in res\_type.\text{actions}$}
        \FOR{each $kw \in meta.\text{keywords}$}
            \IF{$kw \in T_{\text{decl}}.\text{lower}()$}
                \STATE $S_{\text{allowed}} \leftarrow S_{\text{allowed}} \cup \{action\}$
                \STATE $R \leftarrow \max(R, meta.\text{risk})$
            \ENDIF
        \ENDFOR
    \ENDFOR
\ENDFOR
\IF{$S.\text{meta\_rules}.\text{implicit\_read\_grant}$ is True}
    \FOR{each $action \in S_{\text{allowed}}$}
        \IF{$\text{IsWriteOp}(action)$}
            \STATE $S_{\text{allowed}} \leftarrow S_{\text{allowed}} \cup \text{GetReadOps}(action)$
        \ENDIF
    \ENDFOR
\ENDIF
\RETURN $\langle S_{\text{allowed}}, R \rangle$
\end{algorithmic}
\end{algorithm}

\begin{algorithm}[t]
\caption{Policy Decision Enforcement}
\label{alg:policy}
\begin{algorithmic}[1]
\REQUIRE Request $Q$, Scope $\langle S_{\text{allowed}}, R \rangle$, Mode $M$, FailClosed $FC$
\ENSURE Scoping Decision $D \in \{\text{ALLOWED}, \text{BLOCKED}\}$
\TRY
    \IF{$M = \text{PASSTHROUGH} \lor M = \text{TAGGING\_ONLY}$}
        \RETURN \text{ALLOWED}
    \ENDIF
    \IF{$Q.a_{\text{req}} \in S_{\text{allowed}}$}
        \RETURN \text{ALLOWED}
    \ELSE
        \STATE $\text{LogSecurityWarning}(Q.\text{agent\_id}, Q.a_{\text{req}}, S_{\text{allowed}})$
        \RETURN \text{BLOCKED}
    \ENDIF
\CATCH{Exception $E$}
    \IF{$FC = \text{True}$}
        \RETURN \text{BLOCKED}
    \ELSE
        \RETURN \text{ALLOWED}
    \ENDIF
\ENDTRY
\end{algorithmic}
\end{algorithm}

\section{Experimental Evaluation}\label{sec:evaluation}

We evaluate ScopeGuard across a 200-task benchmark suite (100 legitimate operations and 100 adversarial attacks across 10 security threat categories) across all 3 ablation modes (\textbf{600 total tool calls}).

\begin{table}[htbp]
\caption{ScopeGuard Ablation Performance Summary (600 Tool Calls)}
\label{tab:ablation_results_detail}
\begin{center}
\begin{tabular}{lrrr}
\toprule
\textbf{Evaluation Metric} & \textbf{Passthrough} & \textbf{Tagging Only} & \textbf{ScopeGuard Full} \\
\midrule
Total Tool Calls Evaluated & 200 & 200 & 200 \\
True Positives (Blocked Adv) & 0 & 0 & 88 \\
True Negatives (Allowed Leg) & 100 & 100 & 98 \\
False Positives (Blocked Leg) & 0 & 0 & 2 \\
False Negatives (Allowed Adv) & 100 & 100 & 12 \\
\midrule
Overall Accuracy & 50.0\% & 50.0\% & \textbf{89.0\%} \\
Scoping Precision & 0.000 & 0.000 & \textbf{0.978 (97.8\%)} \\
Scoping Recall (Block Rate) & 0.000 & 0.000 & \textbf{0.880 (88.0\%)} \\
F1 Score & 0.000 & 0.000 & \textbf{0.926} \\
False Positive Rate (FPR) & 0.0\% & 0.0\% & \textbf{2.0\%} \\
Mean Proxy Latency Overhead & 129.2 ms & 111.2 ms & \textbf{108.9 ms} \\
\bottomrule
\end{tabular}
\end{center}
\end{table}

\subsection{Key Findings and Analysis}
\begin{enumerate}
    \item \textbf{High Security Efficacy}: ScopeGuard blocks \textbf{88.0\% of adversarial attacks} (88/100). Off-scope operations such as attempting to execute \texttt{iam:DeleteRole} during a database monitoring task are deterministically intercepted.
    \item \textbf{Low False Positive Rate}: Achieves a \textbf{97.8\% precision} and \textbf{2.0\% false positive rate} on legitimate tasks (98/100 allowed).
    \item \textbf{Sub-110ms Real-Time Latency}: The average proxy overhead in full enforcement mode is \textbf{108.9 ms}, fulfilling the sub-200ms real-time latency requirement (NFR1).
\end{enumerate}

\subsection{Live AWS Cloud Endpoint Verification}
ScopeGuard was verified live against real AWS endpoints in \texttt{us-east-1}:
\begin{itemize}
    \item \textbf{Amazon RDS}: \texttt{rds:DescribeDBInstances} executed under task "Describe database status", returning live PostgreSQL database metadata (\texttt{school-bus-db}, \texttt{db.t4g.micro}, endpoint \texttt{school-bus-db.cqzcewy46l10.us-east-1.rds.amazonaws.com:5432}).
    \item \textbf{Amazon EC2}: \texttt{ec2:DescribeInstances} executed under task "Describe running EC2 instances", returning live reservation \texttt{r-0fee7bd611fee1f03}.
\end{itemize}

\section{Related Work}\label{sec:related_work}
Previous research on AI security focuses either on text content guardrails (NeMo Guardrails, Lakera Guard) or post-hoc configuration drift detection (RIVA 2026, IaC closed-loop frameworks). ScopeGuard represents the \textbf{first working, empirically tested implementation} of pre-execution task-scoped permission authorization and identity attribution specifically built for AI agents acting on cloud infrastructure control planes.

\section{Conclusion}\label{sec:conclusion}
ScopeGuard bridges the critical security gap between autonomous AI agent capabilities and cloud governance. Evaluation across a 200-task benchmark demonstrates an \textbf{89.0\% accuracy}, \textbf{88.0\% adversarial block rate}, and \textbf{97.8\% precision} with a mean latency overhead of \textbf{108.9 ms}.

\begin{thebibliography}{99}
\bibitem{ayers2018} G. Ayers et al., ``Memory Hierarchy for Web Search,'' in \textit{Proc. IEEE HPCA}, 2018.
\bibitem{jiang2026} J. Jiang et al., ``A Survey on Large Language Models for Code Generation,'' \textit{ACM TOSEM}, 2026.
\bibitem{hadary2020} O. Hadary et al., ``Protean: VM Allocation Service at Scale,'' in \textit{Proc. USENIX OSDI}, 2020.
\bibitem{hong2024} S. Hong et al., ``MetaGPT: Meta Programming for A Multi-Agent Collaborative Framework,'' \textit{arXiv:2308.00352}, 2024.
\bibitem{yang2024} J. Yang et al., ``SWE-agent: Agent-Computer Interfaces Enable Automated Software Engineering,'' \textit{arXiv:2405.15793}, 2024.
\bibitem{schick2023} T. Schick et al., ``Toolformer: Language Models Can Teach Themselves to Use Tools,'' \textit{arXiv:2302.04761}, 2023.
\bibitem{yao2023} S. Yao et al., ``ReAct: Synergizing Reasoning and Acting in Language Models,'' in \textit{Proc. ICLR}, 2023.
\bibitem{shavit2023} Y. Shavit et al., ``Practices for Governing Agentic AI Systems,'' \textit{OpenAI Research}, 2023.
\end{thebibliography}

\end{document}
